\section{Safety Backlash 与过渡期：拒答边界本身就是对齐选择}

算法和数据之外，开放对齐还面对价值冲突。Llama 2 Chat 的安全策略引发一部分用户反弹，随后出现“uncensored” fine-tunes；同一时期，大量 UltraChat、OpenChat、XwinLM、OpenHermes 等模型用不同数据和 recipe 快速迭代。课程把这段时期视为 expectations 尚未稳定的过渡带。

\subsection{Llama 2 safety backlash}

前一章的“方法—模型—评测”闭环默认大家对好行为已有共识，但安全 backlash 恰好说明这个前提并不成立。本节要回答的不是模型是否需要安全约束，而是约束边界由谁定义、误拒成本由谁承担。模型需要在危险请求、合法敏感问题、幽默和虚构之间划界；如果数据过度奖励拒答，模型会在无害请求上也退缩。读图时不要把社区反弹理解成“安全无用”，而要看到安全行为必须结合用户场景、政策、可控性和误拒成本评估。

\lecturefigure{slide-046.jpg}{Llama 2 Chat 的 safety backlash：chat model 应多安全}{official deck, p.~46}

\teachervoice{00:28:55--00:29:52，Nathan 把 Llama 2 backlash 视为 defining moment：模型的拒答风格让社区意识到“safe”不是单一数值，而是用户与开发者对可控性的分歧。}

\subsection{“Uncensored” models 的目标和风险}

Llama 2 backlash 之后，社区很快给出一个可操作但粗糙的反方向：从训练数据中移除显式拒答。这个选择把“用户控制”放在更高位置，却没有自动提供新的危害判断机制。一些 fine-tune 通过删除 ``Sorry''、``as a language model'' 等样本追求不拒答；读图时应比较的是 refusal style 与真实 harmful compliance，而不是只数道歉短语。它们证明训练数据能快速改变 refusal distribution，也暴露“去掉表面拒答”不等于建立更好的安全推理；模型可能只学会更少说安全模板。

\lecturefigure{slide-047.jpg}{“Uncensored” models：删除拒答样本以提高用户控制}{official deck, p.~47}

\begin{warningbox}{拒答率不是完整 safety metric}
需要至少区分 harmful compliance、benign refusal、事实错误、隐私泄露和可逆性。低拒答率可能提高合法任务可用性，也可能增加危险帮助；高拒答率可能更安全，也可能只是模板化回避。应按风险类别和用户 intent 分层评估。
\end{warningbox}

\teachervoice{00:29:52--00:30:54，讲者说明 uncensored 模型通过移除拒答式训练样本改变 behavior，这条路线延续至今。课堂没有把它当作安全答案，而把它作为社区价值选择的证据。}

\subsection{过渡期模型：配方扩散但缺少统一信号}

随着 instruction data 与训练代码扩散，模型数量快速增长，但很难判断哪个 recipe 真正更好。UltraChat、OpenChat、XwinLM、OpenHermes 等把合成数据、真实对话、不同 base model 和偏好优化混合；模型差异逐渐转化为评测问题。

\lecturefigure{slide-048.jpg}{UltraChat、OpenChat、XwinLM、OpenHermes 等过渡期模型}{official deck, p.~48}

\begin{knowledgebox}{过渡期的真正产物}
这段时期留下的不只是模型，而是数据集、chat template、训练框架和评测入口。很多模型很快过时，基础设施却成为后来 Zephyr、Tulu 和现代 post-training recipe 的组件。判断历史贡献时应看可复用 artifact，而不只看发布当天的排名。
\end{knowledgebox}

\subsection{本章小结}

Safety backlash 说明 alignment 是对输出分布和拒答边界的选择；过渡期模型说明当配方快速扩散时，社区最缺的是可信反馈。下一章进入四套评测工具，比较人类偏好、LLM-as-a-judge 与静态 benchmark 各自适合开发还是发布决策。

